\documentclass[10pt]{article}

% ===== Cheatsheet layout (2 pages target) =====
\usepackage[a4paper,margin=0.45cm]{geometry}
\usepackage{multicol}
\setlength{\columnsep}{0.35cm}
\setlength{\parindent}{0pt}
\setlength{\parskip}{0pt}
\usepackage{microtype}
\pagestyle{empty}

% Multicol tightening
\setlength{\premulticols}{0pt}
\setlength{\postmulticols}{0pt}
\setlength{\multicolsep}{0pt}

% ===== Minted formatting =====
\usepackage{minted}
\usepackage{etoolbox}
\usepackage{relsize}

% Tight box padding/thickness (square boxes)
\setlength{\fboxsep}{0.8pt}
\setlength{\fboxrule}{0.25pt}

\setminted{
  fontsize=\relsize{-1.7},,
  baselinestretch=0.90,
  breaklines,
  breakanywhere,
  breaksymbolleft=,
  tabsize=2,
  autogobble,
  frame=single,
  framesep=0.8pt,
  xleftmargin=0em,
  linenos=false
}

% Tighten vertical space between code blocks (careful but safe)
\BeforeBeginEnvironment{minted}{\vspace{-0.60em}}
\AfterEndEnvironment{minted}{\vspace{-0.75em}}

\linespread{0.92}

\begin{document}

\begin{multicols*}{3}
\raggedcolumns

% ===================== VIM =====================
\begin{minted}{javascript}
//BEFORE STARTING, open a terminal in vim and run
  :set splitright
  :set splitbelow
//MOVEMENT
  //Move to beginning of line: Press 0
//Searching for words
  /word (or prefix of a word) + ENTER
//UNDO: Press U
//REDO: Ctrl + r (undo the undo)
//CUT
  //Cut current line: D D
  //Cut from cursor to end of line: Shift + D
//Replace
  :%s/old/new/g //find old and replace with new everywhere in file
  :%s/old/new/gc //find old, ask for confirmation b4 each replacement
//Commands
  :w //Save, will show compile errors after that
  :q! //quit without saving
  :wq //Save and quit
//MULTIPLE TERMINALS
  :new filename //open another terminal above current terminal
  :vnew filename //open another terminal on the right of current terminal
  :only //close all the windows except the one you are in
//YANK
  yy //yank current line
  3yy //yank 3 lines from cursor onwards
  y$ //yank from cursor to end of line
  y0 //yank from cursor to start of line
//Yanking in visual mode
  1. 0, move cursor to start of line
  2. Press V to enter VISUAL mode
  3. Use j (down) or k (up) to select lines
  4, Press y, yank and auto go back to NORMAL
  //can also do the same thing with cut, press d instead of y
//can copy and paste across terminals opened in the same vim terminal
//PASTE: Press P (paste after cursor)
\end{minted}
% ===================== Basics =====================
\begin{minted}{javascript}
counters = new Counter[3]; //instantiating it
for (int i = 0; i < counters.length; i++) {
    counters[i] = new Counter(i);
} //arr.length
int[] num = {10,20,30};
int[] num = new int[] {10,20,30};
//int[] auto fill with 0, boolean[] auto fill with false
for (Object_type obj: array_name) {...}
//running through array with no indexing
String.format("...",...);
%d (numbers), %.3f (floats, .3 is 3dp), %s (strings), %b (boolean)
System.out.println("...");
if (condition) {}
else {}
while (condition) {}
for( init ; condition ; interation)
//break and continue still works
\end{minted}
% ===================== Compiler =====================
\begin{minted}{javascript}
javac *.java //to compile everything
java Main < inputs/ex1.3.in > OUT
diff OUT outputs/ex1.3.out
vim -d OUT outputs/ex1.3.out
\end{minted}
% ===================== 3) Encapsulation =====================
\begin{minted}{javascript}
//when created, class fields assigned default values
  //whole numbers assigned 0, real numbers 0.0, boolean false, char \u0000, object reference null
\end{minted}
% ===================== 5) Constructors =====================
\begin{minted}{javascript}
//this keyword
  //ALWAYS use it when using fields in a class
\end{minted}
% ===================== 7) Class Fields =====================
\begin{minted}{javascript}
public static int x = 100; //if you reassign will reassign for everything
//final keyword: does not allow the value of the field to be changed
//use public static final for constant values
\end{minted}
% ===================== 8) Class Methods =====================
\begin{minted}{javascript}
class Circle{ //static method
  public static int getCircleCount() {...}}
//CANNOT use instance fields in static method
  //can use class fields in instance method
public static final void main(String[] args) {...}
\end{minted}
% ===================== 11) Inheritence =====================
\begin{minted}{javascript}
class ColoredCircle extends Circle {
  private Color color;
  public ColoredCircle(Point center, double radius, Color color) {
    super(center, radius);
    this.color = color;
  }}
//ONLY public fields and methods in parent will be in child, including public class fields/methods
Circle c = new ColoredCircle(p, 0, blue); //widening type conversion
ColoredCircle cc = c; //compile error, compiler treats c as Circle, dk that its RTT is ColouredCircle
ColoredCircle cc = (ColoredCircle) c; // narrowing type conversion
//if child extends Parent with static method, you can call inherited static method using the child className too
\end{minted}
% ===================== 12) Overriding =====================
\begin{minted}{javascript}
//method signature: name, type, order and number of parameters
  class C {
    A foo(B1 x, B2 y) {...}} // C::foo(B1, B2)
//method descriptor: method signature + return type
  //A C::foo(B1, B2)
//overriding: must have same method descriptor as instance method in parent class
  //return type may be a subtype of the og return type
  //parameters MUST be exactly the same
  //dc about parameter names
  //CANNOT override class method
class ColouredCircle() {
  @Override
  public String toString() {...}
}
//toString method, String Object::toString()
  //called automatically during string concatenation using the + operator
  String s = "Circle c is" + c; //Circle c is Circle@1ce92674
  //overriden toString will still be called automatically
super.methodName();//call overriden method
\end{minted}
% ===================== 13) Overloading =====================
\begin{minted}{javascript}
//Overloading: class has access (defined or inherited) to 2 or more methods with SAME name but different method signatures
  //order, number or type of parameters different, dc about names
  //if only difference is names or no difference at all between 2 methods in SAME class, compile error
  //return type can be different, CANNOT differ only by return type
class Circle {
  public boolean contains(Point p) {...}
  public boolean contains(Point p, double y){...}
}
  class Circle {
    public Circle(Point c, double r) {...}
    public Circle() { this(new Point(0, 0), 1); }
  }
//CAN overload static class methods
  //can also overload inherited public class methods
\end{minted}
% ===================== 14) Polymorphism =====================
\begin{minted}{javascript}
class Circle {
  ...
  @Override
  public boolean equals(Object obj) {
    if (obj instanceof Circle) {
      Circle circle = (Circle) obj;
      return (circle.c.equals(this.c) && circle.r == this.r);
    }
    return false;
  }} //== DOES NOT call equals method in java
//obj instanceof operator
  //check if runtime type of obj is a subtype of Circle, obj<:Circle
\end{minted}
% ===================== 16) Liskov substitution principle (LSP) =====================
\begin{minted}{javascript}
//final keyword: prevent class from being inherited
  final class Circle {...}
  //or can allow inheritence but prevent a specific method from being overriden
  class Circle { public final boolean contains(Point p) {...} }
\end{minted}
% ===================== 17) Abstract class =====================
\begin{minted}{javascript}
//Abstract class (is-a relationship)
  //class that cannot be instantiated
  //if instantiated, will produce an error
abstract class GraphicObject {
  public int x;
  public int y;
  GraphicObject(int x, int y) {
      this.x = x;
      this.y = y;
  }
  void moveTo(int newX, int newY) {
      this.x = newX;
      this.y = newY;
  } //concrete method
  abstract String draw(int x);
  abstract void resize(); //no method body
}
class Circle extends GraphicObject {
  Circle(int x, int y) {
      super(x, y);
  }
  @Override
  String draw(int level) {
      System.out.println("Drawing a circle at " + x + "," + y);
  } //must have same param and return type
  @Override
  void resize() {
      System.out.println("Resizing the circle.");
  }
}
//subclass either implement ALL abstract methods or be declared abstract as well
  //class that inherits must implement ALL abstract methods in parent and grandparent
//NOT all methods have to be abstract, can have no abstract method
//class with 1 abstract method MUST be declared abstract
//can contain class methods, call with className
\end{minted}
% ===================== 18) Interface =====================
\begin{minted}{javascript}
//Interface (can-do relationship)
interface GetAreable {
  double getArea();
} //all methods in interface declared public abstract by default
class Flat extends RealEstate implements GetAreable {
  private int numOfRooms;
  @Override
  public double getArea() {...}
}
//concrete class MUST override ALL abstract methods from interface
//abstract class dont need
class SmartPhone implements Camera, Phone, WebBrowser {...}
//can also extend from 1 or more interfaces
  //interface cannot inherit from class
  //interface can implement other interfaces
    //if both have same concrete method, compile error
interface GetAreable {
  double getArea();
  default double getAreaInSquareFeet() {
    return getArea() * 10.7639;
  }
}
//can have static methods, call with interface name, no override
//can have private methods, only callable within interface
//ALL variables in interface public static final only, constants
//NO constructors
//NOT allowed to override toString() and equals()
//if extends from class and implements interface and both have same method, super will always call from class
\end{minted}
% ===================== 19) Wrapper class =====================
\begin{minted}{javascript}
//Wrapper class: class version of a primitive type
Integer b = Integer.valueOf(5);
Integer x = new Integer(5);  //use is discouraged
int j = b.intValue(); //use this to access the value
//auto boxing and unboxing
Integer a = 5; //autoboxing: int -> Integer
int j = i; //unboxing: Integer -> int
//unboxing can throw NullPointerException if object is null
Double sum = 0.0;
for (int i = 0; i < Integer.MAX_VALUE; i++) {
  sum += i;
}
//should always use equals method when comparing wrapper objects, override so it compares the values not reference type
Double d = 4 //compile error
//cannot simultaneously widen int to Double and autobox double to Double
  //can widen then box or box then widen
\end{minted}
% ===================== 21) Variance =====================
\begin{minted}{javascript}
//variance of types
//covariant if S<:T implies C(S) <: C(T)
//contravariant if S<:T implies C(T) <: C(s)
//invariant if its not covariant or contravariant
//arrays are covariant
  //S<:T, then S[]<:T[]
Integer[] intArray = new Integer[2] { Integer.valueOf(10), Integer.valueOf(20) };
Object[] objArray;
objArray = intArray;
objArray[0] = "Hello!"; // compiles but RT error
\end{minted}
% ===================== 22) Exceptions =====================
\begin{minted}{javascript}
try { ... }
catch (FileNotFoundException e) { ... }
catch (InputMismatchException e) { ... }
catch (NoSuchElementException e) { ... }
finally { ... }
//try blocks: runs first, stops running once exception is thrown
//catch blocks are checked in order
  //first catch block that has an execution type compatible with the type of the thrown exception (type or subtype) will be selected
  //if ExceptionX <: ExceptionY
  catch(ExceptionY e) { ... } 
  catch(ExceptionX e) { ... }//compile error
  //bubble up, runs through catch blocks in method first before being thrown to caller
  //if error caught inside method, code after try and catch block will still run
//finally block (optional): runs after all the try and catch blocks, usually for cleanup tasks, ALWAYS RUNS
class Circle {
  ... 
  public Circle(Point c, double r) throws IllegalArgumentException {
    if (r < 0) {
      throw new IllegalArgumentException("radius cannot be negative.");
    }
    ...
  }
} //throw causes the method to return immediately
//unchecked exceptions: subclass of RuntimeException, cause runtime errors
//checked exceptions: subclass of Exception
class Main {
  static FileReader openFile(String filename) { return new FileReader(filename); }
  public static void main(String[] args) { openFile(); }
} //wont compile cause FileNotFoundException is not handled
//checked exceptions are enforced by compiler
  //if method can throw it, must either declare it with throws or catch it inside the method
  //callers must either catch or declare it too
class Main {
  static FileReader openFile(String filename) {
    try { return new FileReader(filename);} 
    catch (FileNotFoundException e) {
      System.err.println("Unable to open " + filename + " " + e); //calls e.toString()
      return null;
    }
  } //since block catches exception, must return something of type FileReader or null
  public static void main(String[] args) { openFile(); }
}
//unchecked exceptions do not need to be declared in throws
  //exception can still happen at runtime but caller is not forced to catch it and no compile error if its ignored
  //code can compile but if theres error will still have a RT error
//Passing the buck: the caller of the method that generates an exception does not need to catch it, can pass it to caller
class Main {
  static FileReader openFile(String filename) throws FileNotFoundException { return new FileReader(filename); }
  public static void main(String[] args) {
    try { openFile(); } 
    catch (FileNotFoundException e) { ... }
  }
}
//creating exceptions: inherit from existing exceptions
class IllegalCircleException extends IllegalArgumentException {
  Point center;
  IllegalCircleException(String message) { super(message); }
  IllegalCircleException(Point c, String message) { ... }
  @Override
  public String toString() { ... } //this will run when trying to get error string
}
//overriden methods are only allowed to throw the same or more specific (subtype) exceptions
//overriden method can also throw nothing, makes the method easier to call
  //code that handles exception will still compile and run, just that it catches nothing
//NOT allowed to throw new or broader checked exceptions
//unchecked exceptions, even if superclass method throws unchecked exception
  //in override can throw nothing or a different exception
  //dont have to delcare that you are throwing an unchecked exception with the throws keyword
    //can still throw unchecked exception in code and still cause code to immediately return
\end{minted}
% ===================== 23) Generics =====================
\begin{minted}{javascript}
//generic types: class or interface that is parameterised by 1 or more type parameters
//type parameter: what is declared in <>, can ONLY be reference types, NEVER primitive types
//type variable: the placeholder for type, declared in definition of generic class or method, S or T in examples
//type arguement: concrete type or type variable passed to a generic type or method when instantiated
//parameterized type: instantiated generic type
class Pair<S,T> {
  private S first;
  private T second;
  public Pair(S first, T second) {
    this.first = first;
    this.second = second;
  }
  public S getFirst() { return this.first; }
  public T getSecond() { return this.second; }
}
//type arguements can be non-generic type, generic type or another type parameter that has been declared (type variable), ONLY reference types
Pair<String,Integer> foo() { return new Pair<String,Integer>("hello", 4); }
Pair<String,Integer> p = foo();
class DictEntry<T> extends Pair<String,T> { ... }
  //if DictEntry<Integer> it will inherit all accessible instance methods and fields from Pair<String, Integer>
class A {
  public static <T> boolean contains(T[] array, T obj) { ... }
} //generic methods
String[] strArray = new String[] { "hello", "world" };
A.<String>contains(strArray, 123); //compile error
//Bounded type parameters
class A {
  public static <T extends GetAreable> T findLargest(T[] array) {
    for (T curr : array) {
      double area = curr.getArea();
      ...
    } //T must be a subtype of GetAreable
    ...
  } //used for both class and interface
} // have to do this as this guarantees that T is able to use the getArea() function, else compile error
class Pair<S extends Comparable<S>,T> implements Comparable<Pair<S,T>> {
  ... //generic interface Comparable<T>
  public Pair(S first, T second) {
    this.first = first;
    this.second = second;
  } //contains method int compareTo(T t) for any concrete class that implements the interface 
  @Override
  public int compareTo(Pair<S,T> s1) {
    return this.first.compareTo(s1.first);
  } //S extends Comparable<S>, self referential bound, comparable to itself
} //implements Comparable<Pair<S,T>>, so compareTo() takes in Pair<S,T> and can be properly overriden
//when you pass in primitive type into generic method, will ALWAYS autobox as type parameters NEVER primitive type
\end{minted}
% ===================== 24) Type Erasure =====================
\begin{minted}{javascript}
class A {
  void foo(Pair<Sring, String> p) { ... }
  void foo(Pair<Integer, Integer> p) { ... }
}
//CANNOT overload when the only difference between the arguements is the type arguements, compile error
//CANNOT used class type variables in static field and method
class D<T> { //static methods must announce its own independent type parameter
  public static <T> T foo(T x) { ... }
}
String result = D.<String>foo("Hello World");
//generics and arrays CANNOT mix
//generic array declaration, T[] with nothing assigned to it is allowed
//overriding generic methods
class Box<T> { void get(T t) { ... } }
class StringBox extends Box<String> {
  @Override
  void get(String s) { ... }
}
\end{minted}
% ===================== 25) Unchecked Warnings =====================
\begin{minted}{javascript}
ArrayList<Pair<String,Integer>> pairList;
pairList = new ArrayList<Pair<String,Integer>>(); // ok
pairList.add(0, new Pair<Double,Boolean>(3.14, true));  // error
ArrayList<Object> objList = pairList;  // error
class Seq<T> {
  private T[] array;
//not allowed to assign Object[] to T[] without casting, thus (T[])
  public Seq(int size) {
//SuppressWarnings cannot apply to an assignment, only declaration
    @SuppressWarnings("unchecked") //also write comment to explain how its safe to use this
    T[] a = (T[]) new Object[size];
    this.array = a;
  } //tells the compiler everything is fine, no need warning
  public void set(int index, T item) {
    this.array[index] = item;
  } //compiler will check type
  public T get(int index) {
    return this.array[index];
  }
}
Queue<Visitor>[] q = (Queue<Visitor>[]) new Queue<?>[10];
//if you do Object[10], will crash at runtime you physically cannot change an Object to a Queue
//unchecked warnings: compiler saying it has done what it can but because of type erasure there could be RT error that it cannot prevent
//Raw types: generic class without <>
void populateSeq(Seq s) {
  s.set(0, 1234);
  s.set("h","e");
}  //can put any type into the class, compiler will not check, NOT allowed in new code (practice)
//only allowed to use when checking object type at RT, obj instanceof Seq
\end{minted}
% ===================== 26) Wildcards =====================
\begin{minted}{javascript}
//wildcard NOT a type variable, CANNOT use in type parameter declaration
  //CAN put inside the bound of type parameter, and methods
//Upper bound wildcard, Seq<? extends T>
public void copyFrom(Seq<? extends T> src) {
  int len = Math.min(this.array.length, src.array.length);
  for (int i = 0; i < len; i++) {
    this.set(i, src.get(i));
  }
}  //generic type, dk exact type but promise its either T or subclass of T
//if method takes in Seq<? extends T>, get must be assigned to T and set must be null if udk whats inside
//Lower bounded wildcard, Seq<? super T>
public void copyTo(Seq<? super T> dest) {
  int len = Math.min(this.array.length, dest.array.length);
  for (int i = 0; i < len; i++) {
      dest.set(i, this.get(i));
  }
} //A that contains T or supertype of T is a subtype of A which contains S or supertype of S
//if method takes in Seq<? super S>, get must be assigned to Object and set must be S if udk whats inside
Seq<? extends Integer> s = new Seq<Integer>(10); //also compiles
//Producer Extends; Consumer Super (PECS)
  //Producer, ? extends T, can safely get T out
  //Consumer, ? super T, can safely put T into it
//<T extends Comparable<? super T>>, T comparable to itself or supertype of itself
//Unbounded wildcards, Seq<?>
void foo(Seq<?> seq) {
  x = seq.get(0);
  seq.set(0, y);
}
//restrictive, x MUST be Object and y has to be null as Seq can be any type
//Generics are INVARIANT with each other, slight exceptions
  //Seq<String> <: Object
  //if MySeq<T> extends Seq<T>, then MySeq<String> <: Seq<String>, also transitive, MySeq<String> subtype of whatever Seq<String> subtype of
Seq<Object> a1 = new Seq<String>(0); 
Seq<Object> a2 = new Seq<Integer>(0); //compile error
Seq<?> a1 = new Seq<String>(0); // Does compile
Seq<?> a2 = new Seq<Integer>(0); // Does compile
//Seq<?> is a sequence of object of some specific but unknown type
//Seq<Object> is a sequence of Object instances with type checking by the compiler
//Seq is a sequence of Object instances without type checking
a instanceof Seq<?> //replace raw types with this
List<?>[] a = new List<?>[10]; //compiles
Comparable<?>[] b = new Comparable<?>[10]; //compiles
//A<?> is a refiable type as no type information is lost during type erasure, thus can create arrays
  //creates an array of A[], can store any A object in it
\end{minted}
% ===================== 27) Type Inference =====================
\begin{minted}{javascript}
Pair<String,Integer> p = new Pair<>();
//java infers that p should be instance of Pair<String, Integer> based on CTT of p
public static <S> boolean contains(Seq<? extends S> seq, S obj) { ... }
A.contains(circleSeq, shape) //can call like non-generic method
\end{minted}

\end{multicols*}

\end{document}